\documentclass[12pt]{article}

\usepackage{geometry}
\usepackage{titling}
\usepackage{amsmath,amssymb}
\usepackage{fontspec}
\usepackage{unicode-math}
\usepackage{array}
\usepackage[dvipsnames]{xcolor}

%\geometry{a4paper, top=2.5cm, bottom=2.5cm, left=2.6cm, right=2.5cm}
%\pretitle{\begin{center}\LARGE\bfseries}
%\posttitle{\end{center}\vspace{-1.5em}}
%\preauthor{\begin{center}\large}
%\postauthor{\end{center}\vspace{-1.0em}}
%\predate{\begin{center}\vspace{-0.8em}\large}
%\postdate{\par\end{center}\vspace{-0.5em}}
\setlength{\droptitle}{-2.8cm}

\setmonofont{DejaVu Sans Mono}
\setmainfont{Libertinus Serif}
\setsansfont{Libertinus Sans}
\setmathfont{STIXTwoMath-Regular.otf}

\newcommand{\Matrix}[1]{\textbf{#1}}
\newcommand{\Vector}[1]{\textbf{#1}}
\newcommand{\MMat}[1]{\mathbf{#1}}
\newcommand{\MVec}[1]{\mathbf{#1}}
\newcommand{\Norm}[1]{\lVert {\mathbf{#1}} \rVert}
\newcommand{\NormBeta}[1]{\lVert \beta \mathbf{#1} \rVert}
\newcommand{\NormConst}[1]{\lVert #1 \rVert}
% \newcommand{\Trans}{\text{\raisebox{0.5ex}{\scalebox{0.75}{T}}}}
\newcommand{\abs}[1]{\lvert {#1} \rvert}
\newcommand{\Avg}[1]{\mathbf{avg}(\MVec{#1})}
\newcommand{\Std}[1]{\mathbf{std}(\MVec{#1})}

\title{Manipal School of Information Sciences\\AME 5151 - Applied Linear Algebra Lab\vspace*{1em}\\Assignment 1\\}

\author{Academic Year: 2026-2027, Semester: 1}

\date{Friday, 11 September, 2026}

\begin{document}
\pagecolor{white!2}

\maketitle

\vspace*{-1em}\noindent{\color{black!70}{\rule{\textwidth}{0.8pt}}}

{\color{black!70}{
    \noindent This assignment maps to the lab course outcome \textbf{CO1}:\vspace*{2pt}\\
    \hspace*{1em}\textbf{Apply Python's legacy libraries for coding vector and matrix-related operations.}
}}

\vspace*{1.5em}
\noindent\hfill {\textbf{Total Marks}: \textbf{5}}
\vspace*{-4pt}\\
\noindent{\color{black!70}{\rule{\textwidth}{0.8pt}}}\\


\noindent Use your implementation of the {\color{blue}{Vector}} class.
Implement the following methods and write unit tests for them
in a separate file named \textbf{test\_vector.py}.

\begin{enumerate}
    \item \textbf{(1 point)} \\
        {\color{blue!80}{\emph{mean}}} method
        that returns the mean of the vector entries.
        The mean is computed as the sum of all entries divided by
        the number of entries.

    \item \textbf{(1 point)} \\
        {\color{blue!80}{\emph{demean}}} method that
        returns a new vector representing the de-meaned version
        of the original vector ({\color{black!80}\textbf{\textit{self}}}).
        The de-meaned vector is obtained by subtracting the mean
        of the original vector from each of its entries.

    \item \textbf{(1 points)} \\
        {\color{blue!80}{\emph{std}}} method
        that returns the standard deviation of the vector entries.

        The standard deviation is computed as the square root of the
        average of the squared deviations from the mean.
        Use the {\color{blue!80}{\emph{demean}}} method to help
        with this computation.

    \item \textbf{(2 point)} \\
         Write meaningful unit tests for your implementation of
         the above methods.
         The unit tests must cover properties of the mean, de-meaned vector,
         and standard deviation mentioned in the lectures.
         You can refer to the lecture notes/slides for guidance on
         what properties to test.

\end{enumerate}

\end{document}